\begingroup
\raggedright
\GaussContentsEntry{\SourceLeaf{661}\noindent\hangindent=1em Allgemeine Tafeln für Aberration und Nutation\GaussContentsLeaders{\mbox{.\ .\ .\ .\ .\ .\ .\ .\ .\ .}}\GaussContentsReference{\mbox{1808 April}\quad\mbox{Seite 123}}}
\GaussContentsEntry{\noindent\hangindent=1em Über eine Aufgabe der sphärischen Astronomie\GaussContentsLeaders{\mbox{.\ .\ .\ .\ .\ .\ .\ .\ .\ .}}\GaussContentsReference{\mbox{1808 Okt.}\quad\mbox{--- 129}}}
\GaussContentsEntry{\noindent\hangindent=1em \textsc{Gauss} an \textsc{von Lindenau} [Zusatz zum vorhergehenden Aufsatz]\GaussContentsLeaders{\mbox{.\ .\ .\ .\ .}}\GaussContentsReference{\mbox{1808 Nov. 30}\quad\mbox{--- 141}}}
\GaussContentsEntry{\noindent\hangindent=1em \textsc{Gauss} an \textsc{von Lindenau} [Bestimmung der Polhöhe]\GaussContentsLeaders{\mbox{.\ .\ .\ .\ .\ .\ .\ .\ .\ .}}\GaussContentsReference{\mbox{1808 Aug. 30}\quad\mbox{--- 146}}}
\GaussContentsEntry{\noindent\hangindent=1em Kurzer Überblick über die Methoden zur Bestimmung der Bahnen der beiden neuen Hauptplaneten\GaussContentsLeaders{\mbox{.\ .\ .\ .\ .\ .\ .\ .\ .\ .\ .\ .}}\GaussContentsReference{\mbox{1809 Sept.}\quad\mbox{--- 148}}}
\GaussContentsEntry{\noindent\hangindent=1em Tafeln zur Mittagskorrektur\GaussContentsLeaders{\mbox{.\ .\ .\ .\ .\ .\ .\ .\ .\ .\ .\ .\ .\ .}}\GaussContentsReference{\mbox{1811 März}\quad\mbox{--- 166}}}
\GaussContentsEntry{\noindent\hangindent=1em Tafel für die Sonnenkoordinaten\GaussContentsLeaders{\mbox{.\ .\ .\ .\ .\ .\ .\ .\ .\ .\ .\ .\ .\ .\ .}}\GaussContentsReference{\mbox{1812 Jan.}\quad\mbox{--- 172}}}
\GaussContentsEntry{\noindent\hangindent=1em Neue Aussicht auf eine Erweiterung des Gebiets der Astronomie\GaussContentsLeaders{\mbox{.\ .\ .\ .\ .}}\GaussContentsReference{\mbox{1813}\quad\mbox{--- 181}}}
\GaussContentsEntry{\noindent\hangindent=1em Refraktionstafeln\GaussContentsLeaders{\mbox{.\ .\ .\ .\ .\ .\ .\ .\ .\ .\ .\ .\ .\ .\ .\ .\ .\ .\ .\ .\ .\ .}}\GaussContentsReference{\mbox{1822}\quad\mbox{--- 185}}}
\GaussContentsEntry{\noindent\hangindent=1em \textsc{Gauss} an \textsc{Schumacher} [Berechnung der wahren Anomalie eines Kometen]\GaussContentsLeaders{\mbox{.}}\GaussContentsReference{\mbox{1843 April 1}\quad\mbox{--- 191}}}
\GaussContentsEntry{\noindent\hangindent=1em \textsc{Gauss} an \textsc{Schumacher} [Grenzen eines Zodiakus]\GaussContentsLeaders{\mbox{.\ .\ .\ .\ .\ .\ .\ .\ .\ .}}\GaussContentsReference{\mbox{1847 Nov. 23}\quad\mbox{--- 194}}}
\par\medskip
\Needspace{5\baselineskip}\noindent\textit{Beobachtungen und Rechnungen.}\par\smallskip\nopagebreak[4]
\GaussContentsEntry{\noindent\hangindent=1em Fortgesetzte Nachrichten über den längst vermuteten neuen Hauptplaneten unseres Sonnensystems [Ceres Elemente I. II. III. IV.]\GaussContentsLeaders{\mbox{.\ .\ .\ .}}\GaussContentsReference{\mbox{1801 Dez.}\quad\mbox{--- 199}}}
\GaussContentsEntry{\noindent\hangindent=1em Fortgesetzte Nachrichten über den neuen Hauptplaneten [Ceres. Wiederauffindung]\GaussContentsLeaders{\mbox{.\ .\ .\ .\ .\ .\ .\ .\ .\ .\ .\ .\ .\ .\ .\ .\ .\ .\ .\ .\ .\ .}}\GaussContentsReference{\mbox{1802 Feb.}\quad\mbox{--- 204}}}
\GaussContentsEntry{\noindent\hangindent=1em Fortgesetzte Nachrichten über den neuen Hauptplaneten. Ceres Ferdinandea [Ceres Elemente V. VI. VII.]\GaussContentsLeaders{\mbox{.\ .\ .\ .\ .\ .\ .\ .\ .\ .\ .\ .\ .\ .}}\GaussContentsReference{\mbox{1802 März}\quad\mbox{--- 205}}}
\GaussContentsEntry{\noindent\hangindent=1em Fortgesetzte Nachrichten. [Ceres. Beobachtungen]\GaussContentsLeaders{\mbox{.\ .\ .\ .\ .\ .\ .\ .\ .}}\GaussContentsReference{\mbox{1802 April}\quad\mbox{--- 209}}}
\GaussContentsEntry{\noindent\hangindent=1em Fortgesetzte Nachrichten. [Ceres. Verglichene Beobachtungen]\GaussContentsLeaders{\mbox{.\ .\ .\ .}}\GaussContentsReference{\mbox{1802 Mai}\quad\mbox{--- 211}}}
\GaussContentsEntry{\noindent\hangindent=1em Fortgesetzte Nachrichten. [Pallas. Elemente I. II.]\GaussContentsLeaders{\mbox{.\ .\ .\ .\ .\ .\ .\ .}}\GaussContentsReference{\mbox{1802 Juni}\quad\mbox{--- 212}}}
\GaussContentsEntry{\noindent\hangindent=1em Fortgesetzte Nachrichten. [Pallas. Elemente III.]\GaussContentsLeaders{\mbox{.\ .\ .\ .\ .\ .\ .\ .\ .}}\GaussContentsReference{\mbox{1802 Juli}\quad\mbox{--- 217}}}
\GaussContentsEntry{\noindent\hangindent=1em Pallas Olbersiana [Vergleich mit den Elementen III.]\GaussContentsLeaders{\mbox{.\ .\ .\ .\ .\ .\ .}}\GaussContentsReference{\mbox{1802 Aug.}\quad\mbox{--- 221}}}
\GaussContentsEntry{\noindent\hangindent=1em \textsc{Gauss} an \textsc{Bode}. [Pallas. Elemente IV.]\GaussContentsLeaders{\mbox{.\ .\ .\ .\ .\ .\ .\ .\ .\ .\ .\ .\ .}}\GaussContentsReference{\mbox{1802 Sept. 10}\quad\mbox{--- 223}}}
\GaussContentsEntry{\noindent\hangindent=1em Ceres Ferdinandea. [Beobachtungen und Vergleiche]\GaussContentsLeaders{\mbox{.\ .\ .\ .\ .\ .}}\GaussContentsReference{\mbox{1802 Juni}\quad\mbox{--- 224}}}
\GaussContentsEntry{\noindent\hangindent=1em Ceres Ferdinandea. [Störungen durch Jupiter. Länge und Radiusvektor. Elemente VIII.]\GaussContentsLeaders{\mbox{.\ .\ .\ .\ .\ .\ .\ .\ .\ .\ .\ .\ .\ .\ .\ .\ .\ .\ .\ .}}\GaussContentsReference{\mbox{1802 Okt.}\quad\mbox{--- 224}}}
\GaussContentsEntry{\noindent\hangindent=1em Pallas Olbersiana. [Ephemeride]\GaussContentsLeaders{\mbox{.\ .\ .\ .\ .\ .\ .\ .\ .\ .\ .\ .\ .\ .\ .\ .}}\GaussContentsReference{\mbox{1802 Okt.}\quad\mbox{--- 226}}}
\GaussContentsEntry{\noindent\hangindent=1em Ceres Ferdinandea. [Störung durch Jupiter. Länge, Radiusvektor und Breite]\GaussContentsReference{\mbox{1802 Nov.}\quad\mbox{--- 227}}}
\GaussContentsEntry{\noindent\hangindent=1em Pallas Olbersiana. [Abstände der Planeten von der Sonne]\GaussContentsLeaders{\mbox{.\ .\ .\ .\ .\ .}}\GaussContentsReference{\mbox{1802 Nov.}\quad\mbox{--- 230}}}
\GaussContentsEntry{\noindent\hangindent=1em Merkurdurchgang vor der Sonne am 3. November 1802\GaussContentsLeaders{\mbox{.\ .\ .\ .\ .}}\GaussContentsReference{\mbox{1802 Dez.}\quad\mbox{--- 231}}}
\GaussContentsEntry{\noindent\hangindent=1em Pallas Olbersiana. [Elemente V. und Vergleiche]\GaussContentsLeaders{\mbox{.\ .\ .\ .\ .\ .\ .\ .}}\GaussContentsReference{\mbox{1802 Dez.}\quad\mbox{--- 232}}}
\GaussContentsEntry{\noindent\hangindent=1em \textsc{Gauss} an \textsc{Bode}. [Pallas. Elemente VI.]\GaussContentsLeaders{\mbox{.\ .\ .\ .\ .\ .\ .\ .\ .\ .\ .\ .\ .}}\GaussContentsReference{\mbox{1803 März 3}\quad\mbox{--- 234}}}
\GaussContentsEntry{\noindent\hangindent=1em Tafeln für die Störungen der Ceres durch Jupiter\GaussContentsLeaders{\mbox{.\ .\ .\ .\ .\ .\ .\ .\ .\ .}}\GaussContentsReference{\mbox{1802 Dez. 26}\quad\mbox{--- 235}}}
\GaussContentsEntry{\noindent\hangindent=1em Pallas. [Beobachtungen]\GaussContentsLeaders{\mbox{.\ .\ .\ .\ .\ .\ .\ .\ .\ .\ .\ .\ .\ .\ .\ .\ .\ .}}\GaussContentsReference{\mbox{1803 April u. Mai}\quad\mbox{--- 244}}}
\GaussContentsEntry{\noindent\hangindent=1em Pallas. [Beobachtungen]\GaussContentsLeaders{\mbox{.\ .\ .\ .\ .\ .\ .\ .\ .\ .\ .\ .\ .\ .\ .\ .\ .\ .\ .}}\GaussContentsReference{\mbox{1803 Juni}\quad\mbox{--- 244}}}
\GaussContentsEntry{\noindent\hangindent=1em Pallas. [Beobachtungen verglichen mit Elementen VI.]\GaussContentsLeaders{\mbox{.\ .\ .\ .\ .\ .\ .}}\GaussContentsReference{\mbox{1803 Juli}\quad\mbox{--- 245}}}
\GaussContentsEntry{\noindent\hangindent=1em Ceres. [Elemente IX.]\GaussContentsLeaders{\mbox{.\ .\ .\ .\ .\ .\ .\ .\ .\ .\ .\ .\ .\ .\ .\ .\ .\ .\ .}}\GaussContentsReference{\mbox{1803 Sept. u. Okt.}\quad\mbox{--- 245}}}
\GaussContentsEntry{\noindent\hangindent=1em Ceres und Pallas. [Pallas Elemente VII.]\GaussContentsLeaders{\mbox{.\ .\ .\ .\ .\ .\ .\ .\ .\ .\ .\ .}}\GaussContentsReference{\mbox{1804 März}\quad\mbox{--- 246}}}
\GaussContentsEntry{\noindent\hangindent=1em Neuer Komet [von 1804. Elemente]\GaussContentsLeaders{\mbox{.\ .\ .\ .\ .\ .\ .\ .\ .\ .\ .\ .\ .\ .\ .}}\GaussContentsReference{\mbox{1804 Mai}\quad\mbox{--- 247}}}
\GaussContentsEntry{\noindent\hangindent=1em \textsc{Gauss} an \textsc{Bode}. [Juno. Elemente I. III.]\GaussContentsLeaders{\mbox{.\ .\ .\ .\ .\ .\ .\ .\ .\ .}}\GaussContentsReference{\mbox{1804 Sept. 30 u. Nov. 12}\quad\mbox{--- 248}}}
\GaussContentsEntry{\noindent\hangindent=1em Über einen Wandelstern. [Juno. Elemente I.)\GaussContentsLeaders{\mbox{.\ .\ .\ .\ .\ .\ .\ .\ .\ .}}\GaussContentsReference{\mbox{1804 Okt.}\quad\mbox{--- 250}}}
\GaussContentsEntry{\noindent\hangindent=1em Einige Nachrichten über den neuen Planeten. [Juno]\GaussContentsLeaders{\mbox{.\ .\ .\ .\ .\ .\ .\ .}}\GaussContentsReference{\mbox{1804 Okt. 6}\quad\mbox{--- 253}}}
\endgroup
